\documentclass[11pt]{amsart}
\usepackage{amssymb}
\usepackage[margin=1.15in]{geometry}
\usepackage[hidelinks]{hyperref}
\hypersetup{pdftitle={Flatness of a standard etale algebra in two variables}}

% Stamp on every page: file name and compilation date.
\newcommand{\stamptime}{\the\numexpr\time/60\relax:\ifnum\numexpr\time-60*(\time/60)\relax<10 0\fi\the\numexpr\time-60*(\time/60)\relax}
\AddToHook{shipout/foreground}{%
 \put(0.5\paperwidth,-\paperheight+0.55cm){\makebox[0pt]{\tiny\ttfamily
  unit/doc/\jobname.tex \ --- \ compiled \today, \stamptime}}}

\newtheorem{theorem}{Theorem}[section]
\newtheorem{proposition}[theorem]{Proposition}
\newtheorem{lemma}[theorem]{Lemma}
\theoremstyle{definition}
\newtheorem{definition}[theorem]{Definition}
\theoremstyle{remark}
\newtheorem{remark}[theorem]{Remark}
\newcommand{\code}[1]{\texttt{#1}}
\newcommand{\mere}[1]{\left\lVert #1\right\rVert}

\title{Flatness of a standard \'etale algebra in two variables}
\date{October 26, 2026}

\begin{document}
\begin{abstract}
Let $A$ be any commutative ring, $F,G\in A[X,Y]$ and $J=F_XG_Y-F_YG_X$.
We show that $A[X,Y][1/J]/(F,G)$ is a flat $A$-module. The proof is
elementary and constructive: flatness is the existence of explanations of
finite relations; the polynomial ring and its localization are flat; a
relation in the quotient is corrected modulo $(F,G)$ using the syzygies
of a regular pair; and the regularity of $F,G$ modulo every finitely
generated ideal of $A$ comes from finite Taylor lifts of the curve
$F=f-T$, $G=g$ and a noetherian argument in a finitely generated subring.
No hypothesis is made on $A$. The argument is due to Claude Quitt\'e.
Everything is checked in Unit.
\end{abstract}
\maketitle

\section{Statement}

Throughout, $A$ is a commutative ring, $R=A[X,Y]$, $F,G\in R$,
$J=F_XG_Y-F_YG_X$ (formal partial derivatives), $S=R[1/J]$, and
\[
 B=S/(F,G).
\]
We write $f,g$ for the images of $F,G$ in $S$, and $u$ for the inverse of
$J$ in $S$. No condition is put on $J$: when $J=0$, $S$ and $B$ are the
zero ring.

\begin{definition}
Let $C$ be an $A$-algebra. A \emph{relation} in $C$ is a family
$a_1,\dots,a_n\in A$, $b_1,\dots,b_n\in C$ with $\sum_i a_ib_i=0$. An
\emph{explanation} of it is a family $c_{ij}\in A$, $y_j\in C$
($1\le j\le m$) with
\[
 b_i=\sum_j c_{ij}y_j\quad(\text{all }i),\qquad
 \sum_i a_ic_{ij}=0\quad(\text{all }j).
\]
$C$ is \emph{flat} over $A$ when every relation merely has an explanation.
\end{definition}

This is the usual equational criterion, and it is equivalent to the
flatness of the $A$-module underlying $C$ in the form used for modules
(one relation explained by a factorisation through a free module).

\begin{theorem}\label{thm:main}
For every commutative ring $A$ and all $F,G\in A[X,Y]$, the algebra
$B=A[X,Y][1/J]/(F,G)$ is flat over $A$.
\end{theorem}

No noetherian hypothesis, decidable equality, or condition on the
characteristic is used.

\section{Explanations, polynomials, localization}

\begin{lemma}\label{lem:ops}
Fix $a_1,\dots,a_n\in A$. The families $(b_i)$ for which the relation is
explained contain $0$ and are closed under sums, under multiplication by
an element of $C$, and under finite sums. If $\varphi:C\to D$ is a map of
$A$-algebras, the image of an explained family is explained.
\end{lemma}

\begin{proof}
For a sum, concatenate the two lists of pairs $(c_{\cdot j},y_j)$; for
$t\cdot b_i$, replace $y_j$ by $ty_j$; for an image, replace $y_j$ by
$\varphi(y_j)$.
\end{proof}

\begin{lemma}\label{lem:poly}
If $C$ is flat over $A$, so is $C[X]$. Hence $A[X,Y]$ is flat over $A$.
\end{lemma}

\begin{proof}
Let $\sum_i a_ig_i=0$ in $C[X]$. Each $g_i$ is merely the class of a list
of coefficients, so there is merely an $L$ such that all coefficients of
all $g_i$ vanish from degree $L$ on. Induction on $L$. If $L=0$ all $g_i$
are $0$. Otherwise write $g_i=c_i+Xt_i$, with $c_i$ the constant term.
The constant term of the relation gives $\sum_ia_ic_i=0$ in $C$, and the
coefficients of $X^{k+1}$ give $\sum_ia_it_i=0$ in $C[X]$, where the
coefficients of the $t_i$ vanish from degree $L-1$ on. So $(c_i)$ is
explained in $C$, hence in $C[X]$, and $(t_i)$ by induction; by
Lemma~\ref{lem:ops}, so is $(c_i+Xt_i)$. Start from $A$, which is flat
over itself (one pair $(b_i,1)$).
\end{proof}

\begin{lemma}\label{lem:loc}
If $C$ is flat over $A$ and $j\in C$, then $C[1/j]$ is flat over $A$.
\end{lemma}

\begin{proof}
Let $\sum_ia_ib_i=0$ in $C[1/j]$. Finitely many fractions merely have a
common denominator: $b_i=p_i/j^d$. Then $\sum_ia_ip_i/j^d=0$, so
$j^e\sum_ia_ip_i=0$ in $C$ for some $e$. The relation
$\sum_ia_i(j^ep_i)=0$ in $C$ is explained; its image in $C[1/j]$,
multiplied by $1/j^{d+e}$, explains $(b_i)$ (Lemma~\ref{lem:ops}). The
element $j$ need not be regular.
\end{proof}

In particular $S=A[X,Y][1/J]$ is flat over $A$.

\section{Correction of a relation}

\begin{lemma}[Syzygies modulo an ideal]\label{lem:syz}
Let $I$ be an ideal of a ring $C$ and $f,g\in C$ with: $fs\in I$ implies
$s\in I$, and $gv\in(f)+I$ implies $v\in(f)+I$. If $fu+gv\in I$, then
merely there is $w$ with $u-gw\in I$ and $v+fw\in I$.
\end{lemma}

\begin{proof}
$gv=(fu+gv)-fu\in(f)+I$, so $v-ft\in I$ for some $t$. Then
$f(u+gt)=(fu+gv)-g(v-ft)\in I$, so $u+gt\in I$. Take $w=-t$.
\end{proof}

\begin{proposition}\label{prop:correction}
Let $C$ be a flat $A$-algebra and $f,g\in C$ such that for every finite
family $a_1,\dots,a_m$ in $A$, with $I=(a_1,\dots,a_m)C$, $f$ is regular
modulo $I$ and $g$ is regular modulo $(f)+I$ in the sense of
Lemma~\ref{lem:syz}. Then $C/(f,g)$ is flat over $A$.
\end{proposition}

\begin{proof}
Let $\sum_ia_ib_i=0$ in $C/(f,g)$ and lift $b_i$ to $p_i\in C$. Then
$\sum_ia_ip_i=fu+gv$ for some $u,v$. With $I=(a_1,\dots,a_n)C$,
$fu+gv\in I$, and Lemma~\ref{lem:syz} gives $w$ with
\[
 u-gw=\sum_ia_iu_i,\qquad v+fw=\sum_ia_iv_i .
\]
Put $p'_i=p_i-fu_i-gv_i$. Each $p'_i$ still lifts $b_i$, and
\[
 \sum_ia_ip'_i=(fu+gv)-f(u-gw)-g(v+fw)=0\quad\text{in }C.
\]
This relation is explained in $C$; its image explains $(b_i)$.
\end{proof}

So Theorem~\ref{thm:main} follows from Lemmas~\ref{lem:poly}
and~\ref{lem:loc} and the following proposition.

\begin{proposition}\label{prop:regular}
For every commutative ring $A$ and $a_1,\dots,a_m\in A$, with
$I=(a_1,\dots,a_m)S$: $f$ is regular modulo $I$, and $g$ is regular
modulo $(f)+I$.
\end{proposition}

\section{Finite Taylor lifts}

Let $E$ be an $A$-algebra and $x_0,y_0\in E$ with $J(x_0,y_0)$
invertible; put $h=F(x_0,y_0)$ and $k=G(x_0,y_0)$.

\begin{lemma}\label{lem:lift}
For every $N$ there are $x(T),y(T)\in E[T]$ with $x(0)=x_0$, $y(0)=y_0$ and
\[
 F(x(T),y(T))\equiv h-T,\qquad G(x(T),y(T))\equiv k\pmod{T^{N+1}}.
\]
Moreover the lift of order $N+1$ is obtained from the lift of order $N$
by adding $aT^{N+1}$ to $x$ and $bT^{N+1}$ to $y$; so the lifts are
compatible modulo $T^{N+1}$.
\end{lemma}

\begin{proof}
Induction on $N$, starting with the constants $x_0,y_0$. At order $N+1$,
changing $x,y$ by $aT^{N+1},bT^{N+1}$ changes the coefficients of
$T^{N+1}$ in $F(x,y)$ and $G(x,y)$ by $F_Xa+F_Yb$ and $G_Xa+G_Yb$ evaluated
at $(x_0,y_0)$, and no lower coefficient. The $2\times2$ system for the two
errors has determinant $J(x_0,y_0)$, a unit.
\end{proof}

\begin{lemma}[Kill]\label{lem:kill}
Let $(f-T)p(T)\equiv0\pmod{T^{N+1}}$ in $E[T]$, where $f=F(x_0,y_0)$, and
$w$ the coefficient of $T^N$ in $p$. Then $p(0)=f^Nw$ and $f^{N+1}w=0$.
If, for a sequence of such data $w_0,w_1,\dots$ (one for each $N$, with
the same $p(0)$), some $w_k$ is a linear combination of
$w_0,\dots,w_{k-1}$, then $p(0)=0$.
\end{lemma}

\begin{proof}
Comparing coefficients, $p_n=fp_{n+1}$ for $n<N$ and $fp_0=0$; this gives
the two equations. If $w_k=\sum_{j<k}\lambda_jw_j$, then
$f^kw_j=f^{k-1-j}f^{j+1}w_j=0$ for $j<k$, so $p(0)=f^kw_k=0$.
\end{proof}

A sequence has such a good initial segment in a noetherian ring in the
constructive sense of Coquand and Persson (the empty list is barred by
the good lists). This applies inside a subring generated by finitely many
elements, which is a quotient of $\mathbb Z[t_1,\dots,t_r]$ and is
noetherian in this sense by Hilbert's basis theorem, provided each $w_N$
is merely given as a polynomial in the generators.

\section{Regularity}

\begin{lemma}[Kernel principle]\label{lem:kernel}
Let $\varphi:R\to E$ be a ring map with $\varphi(J)$ invertible and
$\varphi(a_i)=0$, and put $x_0=\varphi(X)$, $y_0=\varphi(Y)$.
\begin{enumerate}
\item If $FP=\sum_ia_iQ_i$ in $R$, then $\varphi(P)=0$.
\item If $\varphi(F)=0$ and $GP=FW+\sum_ia_iQ_i$ in $R$, then
 $\varphi(P)=0$.
\end{enumerate}
\end{lemma}

\begin{proof}
(1) Let $E'\subseteq E$ be the subring generated by $x_0,y_0$, the inverse
of $J(x_0,y_0)$, and the images of the coefficients of $F,G,P$. It is
finitely generated, every element of $E'$ is merely a polynomial in the
generators, and the inclusion $E'\to E$ is injective. Run the lifts of
Lemma~\ref{lem:lift} in $E'$, for $F,G$ with coefficients in $E'$. For any
$x(T),y(T)\in E[T]$, the map $R\to E[T]$, $X\mapsto x(T)$, $Y\mapsto y(T)$,
sends each $a_i$ to $0$, so $F(x,y)P(x,y)=0$ in $E[T]$; since $E'\to E$ is
injective, the same holds in $E'[T]$ for the lifts. Thus
$(f-T)P(x(T),y(T))\equiv0\pmod{T^{N+1}}$ for every $N$, with
$P(x(T),y(T))$ of constant term $P(x_0,y_0)=\varphi(P)$, and
Lemma~\ref{lem:kill} gives $\varphi(P)=0$ in $E'$, hence in $E$.

(2) The same, with the roles of $F$ and $G$ exchanged (Jacobian $-J$):
since $F(x_0,y_0)=0$, the lift of Lemma~\ref{lem:lift} gives
$G\equiv g-T$ and $F\equiv0$, so the identity becomes
$(g-T)P(x,y)\equiv F(x,y)W(x,y)\equiv0$ modulo $T^{N+1}$.
\end{proof}

\begin{proof}[Proof of Proposition~\ref{prop:regular}]
\emph{$f$ modulo $I$.} Let $fs\in I$. Write $s=p/J^d$ with $p\in R$; then
$J^dfs$, the image of $Fp$, lies in $I$. An element of $R$ whose image lies
in $I$ satisfies $J^kr=\sum_ia_iQ_i$ in $R$ for some $k,Q_i$: write the
coefficients with a common denominator and clear it with the kernel of
the localization. So $F(J^kp)=\sum_ia_iQ_i$. Apply
Lemma~\ref{lem:kernel}(1) to $E=S/I$: $J^kp=0$ in $S/I$, so $p=0$ there, and
$s\in I$.

\emph{$g$ modulo $(f)+I$.} Let $gv-ft\in I$ and $v=p/J^d$. Then the image of
$Gp$ lies in $(f)+I$, so $J^kGp=FW+\sum_ia_iQ_i$ in $R$. Apply
Lemma~\ref{lem:kernel}(2) to $E=S/(f)+I$ (where $F$ vanishes): $J^kp=0$, so
$p=fc+\sum_ia_ic_i$ in $S$, and $v-f(u^dc)=u^d(p-fc)\in I$.
\end{proof}

This completes the proof of Theorem~\ref{thm:main}. \qed

\begin{remark}
The argument also gives, for every finitely generated ideal $I$ of $A$,
$IS\cap(F,G)S=F\,IS+G\,IS$.
\end{remark}

\section{Formalisation}

Everything above is checked in Unit, in \code{unit/sag/etale-flat}, over
univalence and replacement (the parameters common to all \code{sag/}
developments), with no other hypothesis:
\begin{center}
\begin{tabular}{ll}
Theorem~\ref{thm:main} & \code{EtaleFlat2.Over.flat}, \code{EtaleFlat2.Over.module-flat}\\
equivalence of the two forms of flatness & \code{FlatRelations.Over.equivalent}\\
Lemma~\ref{lem:ops} & \code{FlatRelations} (\code{Ops}, \code{Image})\\
Lemma~\ref{lem:poly} & \code{PolynomialFlat}\\
Lemma~\ref{lem:loc} & \code{LocalizationFlat}\\
Lemma~\ref{lem:syz} & \code{PairSyzygy}\\
Proposition~\ref{prop:correction} & \code{QuotientFlat}\\
Lemma~\ref{lem:lift} & \code{TaylorTargets}, \code{TaylorCompatibility}\\
Lemma~\ref{lem:kill} & \code{StableAnnihilator.Kernel}, \code{CodedSubring.Vanish}\\
Lemma~\ref{lem:kernel} & \code{RegularPolynomial}, \code{RegularCoded}\\
Proposition~\ref{prop:regular} & \code{RegularPair}, \code{LocalizedReduction}
\end{tabular}
\end{center}
Hilbert's basis theorem in the sense of Coquand--Persson is
\code{sag/noetherian}; the coded subring and Lemma~\ref{lem:kill} along a
merely coded sequence are \code{sag/jacobian/CodedSubring}, following
\code{doc/jacobian-noetherian.pdf}. Equalities are equalities in the actual
rings (coefficient quotients, fractions, quotient rings), not of polynomial
functions. The check: \code{./segc sag/etale-flat all}, 121 modules,
6742 declarations, no failure, about 2 seconds.

\end{document}
